\documentclass[french]{article}

\usepackage[utf8]{inputenc}
\usepackage[T1]{fontenc}
\usepackage{lmodern}
\usepackage{geometry}
\usepackage{graphicx}
\usepackage{babel}
\usepackage{amsmath}
\usepackage{amsthm}
\usepackage{amssymb}
\usepackage{hyperref}
\usepackage{stmaryrd}
\usepackage{enumitem}
\usepackage[most]{tcolorbox}
\usepackage{dsfont}
\usepackage{multirow}
\usepackage{etoc}
\usepackage{titlesec}
\usepackage{esint}
\usepackage{cancel}
\usepackage{mathdots}
\usepackage{indentfirst}

\setlength{\parindent}{0pt}

\setlist[itemize]{label=$\bullet$}
\setlist{before=\leavevmode, leftmargin=*}

\renewcommand{\P}{\mathbb{P}}
\newcommand{\N}{\mathbb{N}}
\newcommand{\Z}{\mathbb{Z}}
\newcommand{\Q}{\mathbb{Q}}
\newcommand{\R}{\mathbb{R}}
\newcommand{\C}{\mathbb{C}}
\newcommand{\K}{\mathbb{K}}
\renewcommand{\L}{\mathbb{L}}
\newcommand{\U}{\mathbb{U}}
\newcommand{\st}{^*}
\newcommand{\tq}{\middle|}
\newcommand{\inv}{^{-1}}
\newcommand{\E}{\mathbb{E}}
\newcommand{\F}{\mathbb{F}}
\newcommand{\V}{\mathbb{V}}
\renewcommand{\ker}{\text{Ker}}
\newcommand{\im}{\text{Im}}
\newcommand{\card}{\text{Card}}
\newcommand{\Tr}{\text{Tr}}
\newcommand{\Vect}{\text{Vect}}
\newcommand{\rg}{\text{rg}}
\newcommand{\vertiii}[1]{{\left\vert\kern-0.25ex\left\vert\kern-0.25ex\left\vert #1 \right\vert\kern-0.25ex\right\vert\kern-0.25ex\right\vert}}
\renewcommand{\o}{\text{o}}
\renewcommand{\O}{\text{O}}
\renewcommand{\d}{\text{d}}
\newcommand{\Id}{\text{Id}}


\theoremstyle{plain}
\newtheorem*{ind}{Indication}

\theoremstyle{definition}

\newtheorem*{ex}{Exemple}
\newtheorem*{rmq}{Remarque}
\newtheorem*{notation}{Notation}
\newtheorem*{rappel}{Rappel}
\newtheorem*{terminologie}{Terminologie}
\newtheorem*{attention}{Attention}
\newtheorem*{conséquence}{Conséquence}

\theoremstyle{remark}
\newtheorem*{app}{Application}

\newtcolorbox[auto counter]{dfn}[1][]{
	enhanced,
	colback=white,
	colframe=black,
	sharp corners,
	boxrule=0.8pt,
	fonttitle=\bfseries,
	colbacktitle=white,
	coltitle=black,
	attach boxed title to top left={yshift=-1mm,xshift=-0.5mm},
	boxed title style={size=minimal, right=2pt, sharp corners, colframe=white},
	title={Définition \thetcbcounter \ifstrempty{#1}{}{ (#1)}},
	before skip=0.5em,
	after skip=0.5em
}

\newtcolorbox[auto counter]{exo}[1][]{
	enhanced,
	arc=1mm,
	colback=white,
	colframe=black,
	coltitle=black,
	fonttitle=\bfseries,
	boxrule=0.8pt,
	shadow={1.2pt}{-1.2pt}{0pt}{black!50},
	attach title to upper={\ },
	title={Exercice \thetcbcounter\ifblank{#1}{}{ #1}},
	before skip=0.5em,
	after skip=0.5em,
	left=2mm, right=2mm, top=1mm, bottom=1mm
}

\newtcolorbox{met}[1][]{
	enhanced,
	arc=1mm,
	colback=white,
	colframe=black,
	coltitle=black,
	fonttitle=\bfseries,
	boxrule=0.8pt,
	shadow={1.2pt}{-1.2pt}{0pt}{black!50},
	attach title to upper={\ },
	title={Point méthode \ifblank{#1}{}{#1}},
	before skip=1em,
	after skip=1em,
	left=2mm, right=2mm, top=1mm, bottom=1mm
}

\newcounter{thmcount}

\newtcolorbox[use counter=thmcount]{thm}[1][]{
	enhanced,
	colback=white,
	colframe=black,
	sharp corners,
	left skip=0mm,
	boxrule=1.3pt,
	fonttitle=\bfseries,
	colbacktitle=white,
	coltitle=black,
	attach boxed title to top left={yshift=-1mm,xshift=-0.5mm},
	boxed title style={size=minimal, right=2pt, sharp corners, colframe=white},
	title={Théorème \thethmcount \ifstrempty{#1}{}{ (#1)}},
	before skip=0.5em,
	after skip=0.5em
}

\newtcolorbox[use counter=thmcount]{prop}[1][]{
	enhanced,
	colback=white,
	colframe=black,
	sharp corners,
	boxrule=0.8pt,
	fonttitle=\bfseries,
	colbacktitle=white,
	coltitle=black,
	attach boxed title to top left={yshift=-1mm,xshift=-0.5mm},
	boxed title style={size=minimal, right=2pt, sharp corners, colframe=white},
	title={Proposition \thethmcount \ifstrempty{#1}{}{ (#1)}},
	before skip=0.5em,
	after skip=0.5em
}

\newtcolorbox[use counter=thmcount]{cor}[1][]{
	enhanced,
	colback=white,
	colframe=black,
	sharp corners,
	boxrule=0.8pt,
	fonttitle=\bfseries,
	colbacktitle=white,
	coltitle=black,
	attach boxed title to top left={yshift=-1mm,xshift=-0.5mm},
	boxed title style={size=minimal, right=2pt, sharp corners, colframe=white},
	title={Corollaire \thethmcount \ifstrempty{#1}{}{ (#1)}},
	before skip=0.5em,
	after skip=0.5em
}

\newtcolorbox[use counter=thmcount]{lem}[1][]{
	enhanced,
	colback=white,
	colframe=black,
	sharp corners,
	boxrule=0.5pt,
	fonttitle=\bfseries,
	colbacktitle=white,
	coltitle=black,
	attach boxed title to top left={yshift=-1mm,xshift=-0.5mm},
	boxed title style={size=minimal, right=2pt, sharp corners, colframe=white},
	title={Lemme \thethmcount \ifstrempty{#1}{}{ (#1)}},
	before skip=0.5em,
	after skip=0.5em
}

\title{Matrices et déterminants}
\date{}

\begin{document}
\tableofcontents

\newpage
\maketitle

Dans tout le chapitre, $\K$ désigne un sur-corps de $\Q$.

\section{Généralités}

Revoir dans le cours de première année :
\begin{itemize}
	\item définitions et opérations sur les matrices ;
	\item structures de $\mathcal{M}_{n,p}(\K)$ et de $\mathcal{M}_n(\K)$ ;
	\item matrices inversibles ;
	\item transposition, matrices symétriques et antisymétriques ;
	\item matrices triangulaires, diagonales, scalaires.
\end{itemize}

[A propos du produit matriciel] "Avec la 3ème main, on écrit le résultat. Corollaire : Dieu n'a pas créé les matrices."

\begin{met}
	\begin{enumerate}
		\item Multiplier une matrice $M$ à gauche par une matrice diagonale revient à multiplier les lignes de $M$ par les scalaires de la diagonale. 
		\item Multiplier une matrice $M$ à droite par une matrice diagonale revient à multiplier les colonnes par les scalaires de la diagonale.
	\end{enumerate}
\end{met}

\section{Représentation matricielle}

\subsection{Représentation de vecteurs}

Soit $E$ un espace vectoriel de dimension $n$ et $\mathcal{B}=(e_1,\ldots,e_n)$ une base de $E$.

\subsubsection*{Représentation d'un vecteur}

A tout vecteur $x=\displaystyle\sum_{i=1}^n x_ie_i$, on associe la matrice colonne :
$$X=\begin{pmatrix}
	x_1\\
	\cdots\\
	x_n
\end{pmatrix}\in\K^n=\mathcal{M}_{n,1}(\K)$$ que l'on appelle \textbf{représentant de $x$ dans la base $\mathcal{B}$}. On la note $M_\mathcal{B}(x)$ ou $\text{Mat}_\mathcal{B}(x)$.

\subsubsection*{Représentation d'une famille de vecteurs}

Si $x_1,\ldots,x_p$ sont des vecteurs de $E$, on appelle \textbf{matrice dans la base $\mathcal{B}$} de la famille $(x_1,\ldots,x_p)$ la matrice $(a_{i,j})_{(i,j)\in\llbracket1,n\rrbracket\times\llbracket1,p\rrbracket}\in\mathcal{M}_{n,p}(\K)$ où :
$$\forall j\in\llbracket1,p\rrbracket,\ x_j=\sum_{i=1}^{n}a_{i,j}e_i.$$
On la note $M_\mathcal{B}(x_1,\ldots,x_p)$ ; ses colonnes sont les représentants de chacun des $x_j$ dans la base $\mathcal{B}$.

\subsection{Matrice d'une application linéaire}

Soit $E$ un espace vectoriel de dimension $p$ et $F$ un espace vectoriel de dimension $n$, respectivement munis de bases $\mathcal{B}=(e_j)$ et $\mathcal{B}'=(f_i)$.\\
Soit $u\in\mathcal{L}(E,F)$. On appelle \textbf{matrice de $u$ dans les bases $\mathcal{B}$ et $\mathcal{B}'$} la matrice de la famille $u(\mathcal{B})$ dans la base $\mathcal{B}'$. On la note $M_{\mathcal{B},\mathcal{B}'}(u)$.\\
La matrice $M_{\mathcal{B},\mathcal{B}'}(u)=(a_{i,j})_{(i,j)\in\llbracket1,n\rrbracket\times\llbracket1,p\rrbracket}$ est caractérisée par les relations :
$$\forall j\in\llbracket1,p\rrbracket,\ u(e_j)=\sum_{i=1}^{n}a_{i,j}f_i.$$
Lorsque $F=E$ et $\mathcal{B}'=\mathcal{B}$, on note $M_\mathcal{B}(u)$ la matrice de l'endomorphisme dans la base $\mathcal{B}$.\\
L'application $u\mapsto M_{\mathcal{B},\mathcal{B}'}(u)$ est un isomorphisme d'espaces vectoriels de $\mathcal{L}(E,F)$ sur $\mathcal{M}_{n,p}(\K)$.\\
En particulier, on a $\dim(\mathcal{L}(E,F))=\dim(\mathcal{M}_{n,p}(\K))=np=\dim(E)\times\dim(F)$.

Soit $A\in\mathcal{M}_{n,p}(\K)$. Il existe une unique application linéaire $u$ de $\K^p$ dans $\K^n$ admettant $A$ pour matrice dans les bases canoniques. On l'appelle \textbf{l'application linéaire canoniquement associée à $A$} (ou \textbf{endomorphisme canoniquement associé à $A$} lorsque $A$ est carrée).

\subsubsection*{Propriétés}

Soit :
\begin{itemize}
	\item $E$ un espace vectoriel de dimension $p$ muni d'une base $\mathcal{B}=(e_j)$ ;
	\item $F$ un espace vectoriel de dimension $n$ muni d'une base $\mathcal{B}'=(f_i)$ ;
	\item $x$ un vecteur de $E$ représenté dans la base $\mathcal{B}$ par la matrice $X$ ;
	\item $y$ un vecteur de $F$ représenté dans la base $\mathcal{B}'$ par la matrice $Y$ ;
	\item $u\in\mathcal{L}(E,F)$ de matrice $A$ dans les bases $\mathcal{B}$ et $\mathcal{B}'$.
\end{itemize}
On a alors l'équivalence :
$$y=u(x)\iff Y=AX.$$
En particulier, l'application linéaire canoniquement associée à la matrice $A\in\mathcal{M}_{n,p}(\K)$ est l'application $X\mapsto AX$.

\begin{exo}
	Soit $M\in\mathcal{M}_{n,p}(\K)$. Montrer que :
	\begin{enumerate}
		\item si $\forall X\in\K^p,\ MX=0$, alors $M=0$ ; %ALCA nulle
		\item si $\forall X\in\K^n,\ X^TM=0$, alors $M=0$ ; %transposer et utiliser la question précédente
		\item si $\forall X\in\K^p,\ \forall Y\in\K^n,\ Y^TMX=0$, alors $M=0$. %utiliser les deux premières questions
	\end{enumerate}
\end{exo}

\begin{exo}
	Soit $M\in\mathcal{M}_n(\K)$.
	\begin{enumerate}
		\item Montrer que si $M$ est antisymétrique, alors $\forall X\in\K^n,\ X^TMX=0$.
		\item Montrer la réciproque en calculant $(X+Y)^TM(X+Y)$ pour $X\in\K^n$ et $Y\in\K^n$.
		\item Montrer que si $M$ est symétrique et si $\forall X\in\K^n,\ X^TMX=0$, alors $M=0$.
	\end{enumerate}
\end{exo}

Soit :
\begin{itemize}
	\item $E$, $F$ et $G$ trois espaces vectoriels de dimensions respectives $r$, $q$ et $p$, munis respectivement de base $\mathcal{B}$, $\mathcal{B}'$ et $\mathcal{B}''$ ;
	\item $u\in\mathcal{L}(E,F)$ de matrice $A=M_{\mathcal{B},\mathcal{B}'}(u)\in\mathcal{M}_{q,r}(\K)$ ;
	\item $v\in\mathcal{L}(F,G)$ de matrice $B=M_{\mathcal{B}',\mathcal{B}''}(v)\in\mathcal{M}_{p,q}(\K)$.
\end{itemize}
Alors $M_{\mathcal{B},\mathcal{B}''}(v\circ u)=BA$.

En particulier, si $F=E$, l'application $u\mapsto M_\mathcal{B}(u)$ est un isomorphisme d'algèbres de $\mathcal{L}(E)$ sur $\mathcal{M}_n(\K)$.

Si $\dim(E)=\dim(F)=n$, la matrice d'une application linéaire $u\in\mathcal{L}(E,F)$ dans des bases $\mathcal{B}$ et $\mathcal{B}'$ de $E$ et $F$ est dans $\mathcal{M}_n(\K)$. Elle est inversible si, et seulement si, $u$ est un isomorphisme, et alors :
$$(M_{\mathcal{B},\mathcal{B}'}(u))\inv=M_{\mathcal{B}',\mathcal{B}(u\inv)}.$$

Exercice Moulin.Mat.08

\subsection{Changement de base}

\subsubsection*{Matrice d'un changement de base}

Soit $\mathcal{B},\mathcal{B}'$ deux bases de $E$. La matrice $P=M_{\mathcal{B}}(\mathcal{B}')$ s'appelle la \textbf{matrice de passage} de $\mathcal{B}$ à $\mathcal{B}'$. Ses colonnes sont, par définition, constituées des composantes des vecteurs de $\mathcal{B}'$ dans la base $\mathcal{B}$.

On remarque que $P=M_{\mathcal{B}',\mathcal{B}}(\text{Id}_E)$, et donc qu'elle est inversible, avec :
$$M_{\mathcal{B}}(\mathcal{B}')\inv=M_{\mathcal{B}'}(\mathcal{B}).$$

Si un vecteur $x$ est représenté dans les bases $\mathcal{B}$ et $\mathcal{B}'$ respectivement par les matrices $X$ et $X'$, on a $X=PX'$.

Réciproquement, si $P$ est une matrice inversible et $\mathcal{B}$ une base de $E$, il existe des bases $\mathcal{B}'$ et $\mathcal{B}''$ de $E$ telles que :
$$P=M_{\mathcal{B}}(\mathcal{B}')=M_{\mathcal{B}''}(\mathcal{B})=M_\mathcal{B}(\mathcal{B}'')\inv.$$

\subsubsection*{Effet d'un changement de bases sur la matrice d'une application linéaire}

\begin{dfn}
	Une matrice $M\in\mathcal{M}_{n,p}(\K)$ \textbf{représente une application linéaire} $u\in\mathcal{L}(E,F)$ s'il existe une base $\mathcal{B}_1$ de $E$ et une base $\mathcal{B}_2$ de $F$ telles que $M_{\mathcal{B}_1,\mathcal{B}_2}(u)$.
\end{dfn}

Soit :
\begin{itemize}
	\item $\mathcal{B}$ et $\mathcal{B}'$ deux bases de $E$, et $P$ la matrice de passage de $\mathcal{B}$ à $\mathcal{B}'$ ;
	\item $\mathcal{C}$ et $\mathcal{C}'$ deux bases de $F$, et $Q$ la matrice de passage de $\mathcal{C}$ à $\mathcal{C}'$ ;
	\item $u\in\mathcal{L}(E,F)$ de matrices $A=M_{\mathcal{B},\mathcal{C}}(u)$ et $A'=M_{\mathcal{B'},\mathcal{C'}}(u)$.
\end{itemize}
Alors $A'=Q\inv AP$.

\begin{dfn}
	Deux matrices $A$ et $B$ de $\mathcal{M}_{n,p}(\K)$ sont \textbf{équivalentes} s'il existe des matrices inversibles $P\in\mathcal{GL}_n(\K)$ et $Q\in\mathcal{GL}_p(\K)$ telles que $B=PAQ$.
\end{dfn}

C'est une relation d'équivalence.

\begin{prop}%prop
	Deux matrices sont équivalentes si, et seulement si, elles représentent la même application linéaire (dans des bases \textit{a priori} différentes, évidemment).
\end{prop}

\subsubsection*{Effet d'un changement de base sur la matrice d'un endomorphisme}

Dans le cas des endomorphismes, on utilise généralement la même base au départ et à l'arrivée.

\begin{dfn}
	Une matrice carrée $M\in\mathcal{M}_n(\K)$ \textbf{représente un endomorphisme} $u\in\mathcal{L}(E)$ s'il existe une base $\mathcal{B}$ de $E$ telle que $M=M_\mathcal{B}(u)$.
\end{dfn}

Soit :
\begin{itemize}
	\item $\mathcal{B}$ et $\mathcal{B}'$ deux bases de $E$, et $P$ la matrice de passage de $\mathcal{B}$ à $\mathcal{B}'$ ;
	\item $u\in\mathcal{L}(E)$ de matrices $A=M_\mathcal{B}(u)$ et $A'=M_{\mathcal{B}'}(u)$.
\end{itemize}
Alors $A'=P\inv AP$.

\begin{dfn}
	Deux matrices $A$ et $B$ de $\mathcal{M}_n(\K)$ sont \textbf{semblables} s'il existe une matrice inversible $P\in\mathcal{GL}_n(\K)$ telle que $B=P\inv AP$.
\end{dfn}

C'est une relation d'équivalence.

\begin{prop}%prop
	Deux matrices sont semblables si, et seulement si, elles représentent le même endomorphisme (dans des bases \textit{a priori} différentes, évidemment).
\end{prop}

\begin{rmq}
	On fera attention à la différence entre les notions d'équivalence et de similitude.
\end{rmq}

\subsection{Dualité, transposition (HP)}

Soit $\mathcal{B}$ une base de $E$ (de dimension finie) et $\mathcal{B}\st$ sa base duale.
\begin{itemize}
	\item Si $Y$ est la matrice (colonne) représentant un élément $\varphi\in E\st$ dans $\mathcal{B}\st$, la matrice (ligne) de la forme linéaire $\varphi$ dans la base $\mathcal{B}$ de $E$ et la base canonique de $\K$ est $Y^T$.
	\item Si $X$ représente un vecteur $x$ de $E$ dans la base $\mathcal{B}$, on a :
	$$\langle\varphi,x\rangle=Y^TX.$$
	\item Si $P$ est la matrice d'un changement de base de $E$, la matrice de changement de bases duales correspondant est $Q=(P^T)\inv$.
\end{itemize}

\begin{exo}[$\star$]
	Soit $E$ et $F$ deux espaces vectoriels de dimension finie. Pour $u\in\mathcal{L}(E,F)$, on définit $u^t\in\mathcal{L}(F\st,E\st)$ par $u^t(\varphi)=\varphi\circ u$.
	\begin{enumerate}
		\item Soit $\mathcal{B}$ une base de $E$ et $\mathcal{C}$ une base de $F$. Montrer que 
		$$M_{\mathcal{C}\st,\mathcal{B}\st}(u^t)=M_{\mathcal{B},\mathcal{C}}(u)^T.$$
		\item Déterminer les noyau et image de $u^t$.
		\item Montrer que toute matrice a même rang que sa transposée.
	\end{enumerate}
\end{exo}

\section{Algèbres $\mathcal{L}(E)$ et $\mathcal{M}_n(\K)$}

\subsection{Matrices inversibles}

\begin{prop}
	Soit $A\in\mathcal{M}_n(\K)$. Les propriétés suivantes sont équivalentes :
	\begin{enumerate}[label=(\roman*)]
		\item $A$ est inversible ;
		\item il existe $B\in\mathcal{M}_n(\K)$ telle que $AB=I_n$ ;
		\item il existe $B\in\mathcal{M}_n(\K)$ telle que $BA=I_n$ ;
		\item pour tout $X\in\K^n$, $AX=0\implies X=0$.
	\end{enumerate}
\end{prop}

\paragraph{Matrices à diagonale dominante} Exercice Moulin.Mat.49

\begin{exo}
	\begin{enumerate}
		\item Montrer que si $A\in\mathcal{M}_n(\K)$ est inversible, alors il existe $P\in\K[X]$ tel que $A\inv=P(A)$.
		\item Peut-on trouver $P\in\K[X]$ tel que $\forall A\in\mathcal{GL}_n(\K)$, on ait $A\inv=P(A)$.
	\end{enumerate}
\end{exo}

\subsection{Centre de $\mathcal{M}_n(\K)$}

\begin{dfn}
	Le \textbf{centre} d'une algèbre (ou plus généralement d'un anneau) $A$ est l'ensemble des éléments de $A$ qui commutent avec tout élément de $A$.
\end{dfn}

\begin{prop}[Caractérisation des homothéties (HP)]%prop
	Soit $E$ un espace vectoriel et $u\in\mathcal{L}(E)$ tel que pour tout $x\in E$, la famille $(x,u(x))$ soit liée. Alors $u$ est une homothétie.
\end{prop}

\begin{exo}
	Montrer que le centre de $\mathcal{L}(E)$ est l'ensemble des homothéties.
\end{exo}

\begin{exo}
	\begin{enumerate}
		\item Déterminer le centre de $\mathcal{M}_n(\K)$. %traduction matricielle de l'exo précédent
		\item Plus précisément, trouver toutes les matrices de $\mathcal{M}_n(\K)$ commutant avec toutes les matrices inversibles.
	\end{enumerate}
\end{exo}

\subsection{Trace}

\subsubsection*{Définitions}

\begin{dfn}
	Si $A=(a_{i,j})\in\mathcal{M}_n(\K)$ est une matrice \textit{carrée}, la \textbf{trace} de $A$ est :
	$$\Tr(A)=\sum_{i=1}^{n}a_{i,i}.$$
\end{dfn}

\begin{prop}%prop
	\begin{enumerate}
		\item La trace est une forme linéaire sur $\mathcal{M}_n(\K)$.
		\item Pour $A\in\mathcal{M}_{n,p}(\K)$ et $B\in\mathcal{M}_{p,n}$, on a :
		$$\Tr(AB)=\sum_{(i,j)\in\llbracket1,n\rrbracket\times\llbracket1,p\rrbracket}a_{i,j}b_{j,i}=\Tr(BA).$$
	\end{enumerate}
\end{prop}

\begin{rmq}
	Dans la formule $\Tr(AB)=\Tr(BA)$, on ne suppose pas $A$ et $B$ carrées.
\end{rmq}

Deux matrices semblables (donc carrées) ont donc même trace, puisque :
$$\Tr(P\inv AP)=\Tr((AP)P\inv)=\Tr(A)$$ ce qui permet de poser la définition suivante.

\begin{dfn}
	Si $u$ est un endomorphisme d'un espace vectoriel $E$ de dimension finie, la \textbf{trace} de $u$ est la trace de sa matrice dans n'importe quelle base de $E$. On la note $\Tr(u)$.
\end{dfn}

\subsubsection*{Trace d'un projecteur}

Soit $p$ la projection sur $F$ parallèlement à $G$, où $F$ et $G$ sont deux sous-espaces vectoriels supplémentaires dans $E$ (toujours supposé de dimension finie). Si $\mathcal{B}_1$ est une base de $F$ et $\mathcal{B}_2$ est une base de $G$, alors $\mathcal{B}$ la concaténation de $\mathcal{B}_1$ et $\mathcal{B}_2$ est une base de $E$, dans laquelle la matrice de $p$ est :
$$M_\mathcal{B}(p)=\begin{pmatrix}
	I_r&0\\
	0&0
\end{pmatrix}\quad\text{avec}\quad r=\dim(F)=\rg(p).$$
On en déduit $\rg(p)=r=\Tr(p)$.

\begin{exo}
	Qu'en est-il pour une symétrie ?
\end{exo}

\begin{met}
	Si on doit prouver qu'une trace est un entier (ou est reliée à des entiers, \textit{via} de la divisibilité par exemple), c'est souvent qu'il y a un projecteur ou une symétrie qui se cache dans le problème.
\end{met}

\subsubsection*{Représentation des formes linéaires sur $\mathcal{M}_n(\K)$ (HP)}

Exercice Moulin.Mat.39 (on en donnera deux démonstrations).

Application : exercice Moulin.Mat.40.

\section{Calcul matriciel et systèmes d'équations linéaires}

\subsection{Rang d'une matrice}

\begin{dfn}
	Le \textbf{rang d'une matrice} $A\in\mathcal{M}_{n,p}(\K)$ est le rang de la famille de ses colonnes dans $\K^n$.
\end{dfn}

Par l'isomorphisme $x\mapsto M_\mathcal{B}(x)$ d'un espace vectoriel $E$ de dimension $n$ muni d'une base $\mathcal{B}$ dans $\K^n$, c'est aussi le rang de toute famille de $p$ vecteurs de $E$ dont la matrice dans $\mathcal{B}$ est $A$.

Le rang d'une matrice $A$ est donc aussi le rang de toute application linéaire représentée par $A$.

Deux matrices équivalentes ont donc même rang. En particulier, la multiplication à gauche ou à droite par une matrice inversible ne change pas le rang.

\begin{prop}%prop
	Si $\rg(A)=r$, alors $A$ est équivalente à $J_r=\begin{pmatrix}
		I_r&0\\
		0&0
	\end{pmatrix}$.
\end{prop}

Exercice Moulin.Mat.32

\begin{cor}%cor
	Le rang d'une matrice est égal au rang de sa transposée, donc aussi au rang de ses lignes.
\end{cor}

\begin{cor}%cor
	Deux matrices $A$ et $B$ de mêmes dimensions sont équivalentes si, et seulement si, elles ont même rang.
\end{cor}

\begin{rmq}
	On en déduit que si $A\in\mathcal{M}_n(\K)$, alors $A$ et $A^T$ sont équivalentes. En fait, dans un sous-corps de $\C$, elles sont mêmes semblables, mais cela est bien plus difficile à démontrer et découle du théorème de réduction de Jordan.
\end{rmq}

\begin{exo}
	Combien y a-t-il de classes d'équivalence :
	\begin{enumerate}
		\item dans $\mathcal{M}_{n,p}(\K)$ pour la relation d'équivalence matricielle ? %$1+\min(n,p)$
		\item dans $\mathcal{M}_n(\K)$ pour la relation de similitude matricielle ? %beaucoup, réduction de Jordan
	\end{enumerate}
\end{exo}

\subsection{Opérations sur les lignes ou les colonnes}

\subsubsection*{Matrices $E_{i,j}$}

Pour $i\in\llbracket1,n\rrbracket$ et $j\in\llbracket1,p\rrbracket$, on note $E_{i,j}$ la matrice de $\mathcal{M}_{n,p}(\K)$ dont tous les termes sont nuls sauf celui de la ligne $i$ et de la colonne $j$ qui vaut $1$. Ces matrices constituent la \textbf{base canonique} de $\mathcal{M}_n(\K)$.

\begin{prop}[Table de multiplication des $E_{i,j}$]%prop
	Lorsque les dimensions sont compatibles, on a :
	$$E_{i,j}E_{k,l}=\delta_{j,k}E_{i,l}.$$
\end{prop}

\begin{ex}
	En prenant $(k,l)=(i,j)$ dans le résultat précédent :
	\begin{itemize}
		\item pour $i\ne j$, la matrice $E_{i,j}$ est nilpotente d'ordre $2$ ;
		\item pour $i=j$, elle est idempotente.
	\end{itemize}
\end{ex}

\subsubsection*{Opérations élémentaires}

\begin{itemize}
	\item Pour $i\ne j$, soit $P_{i,j}=I_n-E_{i,i}-E_{j,j}+E_{i,j}+E_{j,i}$ (\textbf{matrice de transposition}). On a $P_{i,j}^2=I_n$.\\
	Échanger dans une matrice $A$ les lignes $i$ et $j$ revient à multiplier $A$ à gauche par $P_{i,j}$.
	\item Pour $\lambda\ne0$, soit $D_i(\lambda)=I_n+(\lambda-1)E_{i,i}$ (\textbf{matrice de dilatation}). On a $D_i(\lambda)D_i(1/\lambda)=I_n$.\\
	Multiplier la ligne $i$ d'une matrice $A$ par $\lambda$ revient à multiplier $A$ à gauche par $D_i(\lambda)$.
	\item Pour $i\ne j$, soit $T_{i,j}(\lambda)=I_n+\lambda E_{i,j}$ (\textbf{matrice de transvection}).
	On a $T_{i,j}(\lambda)T_{i,j}(-\lambda)=I_n$.\\
	Ajouter à la ligne $i$ d'une matrice la ligne $j$ multipliée par $\lambda$ revient à multiplier $A$ à gauche par $T_{i,j}(\lambda)$.
\end{itemize}

Toutes ces opérations conservent le rang de $A$ car les matrices $P_{i,j}$, $D_i(\lambda)$ et $T_{i,j}(\lambda)$ sont inversibles.

Par transposition, on déduit les opérations analogues sur les colonnes.

Rappelons (cours de première année) que toute matrice inversible de $\mathcal{M}_n(\K)$ peut être transformée en $I_n$ par une suite d'opérations élémentaires sur les lignes (méthode du pivot de Gauss).

\begin{exo}
	\begin{enumerate}
		\item Montrer que $\mathcal{GL}_n(\K)$ est engendré par les matrices de transvection, de dilatation et de transposition.
		\item Montrer que l'on peut se passer des transpositions.
		\item Est-ce que $\mathcal{GL}_n(\K)$ est engendré par les transvections ? Sinon, quel est le sous-groupe de $\mathcal{GL}_n(\K)$engendré par ces dernières ?
	\end{enumerate}
\end{exo}

\subsection{Application aux systèmes linéaires}

Soit $A\in\mathcal{M}_{n,p}(\K)$.

\subsubsection*{Calcul du rang}

\begin{itemize}
	\item Si la première colonne de $A$ est nulle, on peut la supprimer sans changer le rang.
	\item Sinon, on peut par des opérations élémentaires sur les lignes transformer $A$ en :
	$$\begin{pmatrix}
		1&*&\cdots&*\\
		0&&&\\
		\vdots&&A'&\\
		0&&&
	\end{pmatrix}.$$
	La première ligne de la matrice obtenue est alors indépendante des autres, ce qui prouve l'égalité $\rg(A)=1+\rg(A')$.
\end{itemize}
On se ramène ainsi ) une matrice de taille strictement inférieure. On termine lorsque $A$ est de rang $0$, c'est-à-dire nulle (ou vide).

\begin{exo}
	Déterminer le rang de la matrice $\begin{pmatrix}
		0&1&4&9\\
		1&0&1&4\\
		4&1&0&1\\
		9&4&1&0
	\end{pmatrix}$.
\end{exo}

\subsubsection*{Inversion d'une matrice}

Si $A$ est inversible et transformée par des opérations élémentaires sur les lignes en $$\begin{pmatrix}
	\lambda&*&\cdots&*\\
	0&&&\\
	\vdots&&A'&\\
	0&&&
\end{pmatrix}$$ avec $\lambda\ne0$, alors $A'$ est inversible.

On peut alors réitérer l'opération sur $A$ et ainsi de suite.

On arrive finalement, par des opérations sur les lignes, à une matrice triangulaire supérieure à coefficients non nuls sur la diagonale.

Réciproquement, si par des opérations sur les lignes d'une matrice $A$, on obtient une telle matrice, alors $A$ est inversible.

\begin{met}
	Pour montrer qu'une matrice $A\in\mathcal{M}_n(\K)$ est inversible et calculer son inverse, on peut résoudre le système $AX=Y$ en le ramenant à un système triangulaire par l'algorithme du pivot de Gauss.\\
	L'expression de $X$ en fonction de $Y$ donne alors la matrice $A\inv$.
\end{met}

\begin{exo}
	Inverser la matrice $\begin{pmatrix}
		1&1&0\\
		1&2&1\\
		1&3&3
	\end{pmatrix}$.
\end{exo}

La méthode précédente donne à chaque étape un système équivalent au système initial. Mais on n'est pas obligé d'utiliser une méthode par équivalence.

\begin{met}
	Si, quel que soit $Y\in\K^n$, le système carré $AX=Y$ entraîne une relation du type $X=\ldots$, alors $A$ est inversible et l'expression de $X$ en fonction de $Y$ donne l'inverse de $A$.
\end{met}

\begin{exo}
	\begin{enumerate}
		\item Inverser la matrice $\begin{pmatrix}
			0&2&3&\cdots&n-1&n\\
			1&0&2&\cdots&n-1&n\\
			1&2&0&\ddots&&\vdots\\
			1&2&3&\ddots&\ddots&n\\
			\vdots&\vdots&\vdots&\ddots&0&n\\
			1&2&3&\cdots&n-1&0
		\end{pmatrix}$.
		\item Retrouver le résultat en introduisant la matrice $D=\text{Diag}(1,2,\ldots,n)$.
	\end{enumerate}
\end{exo}

\section{Écriture par blocs}

\subsection{Matrices par blocs}

Il est parfois utile d'écrire une matrice en mettant en évidence certaines de ses sous-matrices.

\begin{ex}
	La matrice $M=\begin{pmatrix}
		1&2&3&4\\
		0&6&8&9\\
		12&7&0&0
	\end{pmatrix}$ peut s'écrire par blocs $M=\begin{pmatrix}
	A&B\\
	C&D
	\end{pmatrix}$ avec :
	$$A=\begin{pmatrix}
		1&2\\
		0&6
	\end{pmatrix}\ ;\ B=\begin{pmatrix}
	3&4\\
	8&9
	\end{pmatrix}\ ;\ C=\begin{pmatrix}
	12&7
	\end{pmatrix}\ \text{et}\ D=\begin{pmatrix}
	0&0
	\end{pmatrix}$$
	mais aussi $M=\begin{pmatrix}
		E\\
		F
	\end{pmatrix}$ avec $E=\begin{pmatrix}
	1&2&3&4\\
	0&6&8&9\\
	\end{pmatrix}$ et $F=\begin{pmatrix}
	12&7&0&0
	\end{pmatrix}$.
\end{ex}

On appelle \textbf{écriture par blocs} toute écriture de la forme $\begin{pmatrix}
	A_{1,1}&\cdots&A_{1,p}\\
	\vdots&&\vdots\\
	A_{n,1}&\cdots&A_{n,p}
\end{pmatrix}$.

\begin{rmq}
	Bien évidemment, pour qu'une telle écriture ait un sens, deux blocs $A_{i,j_1}$ et $A_{i,j_2}$ ont le même nombre de lignes, et deux blocs $A_{i_1,j}$ et $A_{i_2,j}$ le même nombre de colonnes.
\end{rmq}

\begin{ex}
	\begin{enumerate}
		\item La matrice $J_r\in\mathcal{M}_{n,p}(\K)$ s'écrit naturellement $J_r=\begin{pmatrix}
			I_r&0_{r,p-r}\\
			0_{n-r,r}&0_{n-r,p-r}
		\end{pmatrix}$. \\
		Dans la pratique, on notera souvent plus simplement $J_r=\begin{pmatrix}
			I_r&0\\
			0&0
		\end{pmatrix}$, en laissant implicites les tailles des trois blocs nuls.
		\item Soit $F$ et $G$ deux sous-espaces vectoriels supplémentaires dans $E$ (de dimension finie). Dans une base $\mathcal{B}$ adaptée à la décomposition $E=F\oplus G$, les matrices de la projection $p$ sur $F$ parallèlement à $G$ et de la symétrie $s$ par rapport à $F$ parallèlement à $G$ s'écrivent par blocs :
		$$M_\mathcal{B}(p)=\begin{pmatrix}
			I_r&0\\
			0&0
		\end{pmatrix}\quad\text{et}\quad M_\mathcal{B}(s)=\begin{pmatrix}
		I_r&0\\
		0&-I_{n-r}
		\end{pmatrix}$$ en notant $n$ la dimension de $E$ et $p$ celle de $F$.
		\item On peut toujours écrire $M\in\mathcal{M}_{n,p}(\K)$ par blocs en utilisant la famille de ses lignes ou ses colonnes :
		$$M=\begin{pmatrix}
			L_1\\
			\vdots\\
			L_n
		\end{pmatrix}\quad\text{et}\quad M=\begin{pmatrix}
		C_1&\cdots&C_p
		\end{pmatrix}.$$
	\end{enumerate}
\end{ex}

\subsection{Opérations sur les matrices par blocs}

\subsubsection*{Transposition par blocs}

Si $A=\begin{pmatrix}
	A_{1,1}&\cdots&A_{1,p}\\
	\vdots&&\vdots\\
	A_{n,1}&\cdots&A_{n,p}
\end{pmatrix}$ alors $A^T=\begin{pmatrix}
A_{1,1}^T&\cdots&A_{n,1}^T\\
\vdots&&\vdots\\
A_{1,p}^T&\cdots&A_{n,p}^T
\end{pmatrix}$.

\subsubsection*{Combinaison linéaire par blocs}

Soit $A$ et $B$ deux matrices de \textit{même taille}, écrites avec des blocs de \textit{tailles compatibles pour l'addition}, alors toute combinaison linéaire de $A$ et $B$ s'écrit par blocs :
$$\lambda A+\mu B=\begin{pmatrix}
	\lambda A_{1,1}+\mu B_{1,1}&\cdots&\lambda A_{1,p}+\mu B_{1,p}\\
	\vdots&&\vdots\\
	\lambda A_{n,1}+\mu B_{n,1}&\cdots&\lambda A_{n,p}+\mu B_{n,p}
\end{pmatrix}.$$

\subsubsection*{Produit par blocs}

\begin{prop}[Produit par blocs]%prop
	Si $A$ et $B$ sont écrites par blocs sous la forme :
	$$A=\begin{pmatrix}
		A_{1,1}&\cdots&A_{1,p}\\
		\vdots&&\vdots\\
		A_{n,1}&\cdots&A_{n,p}
	\end{pmatrix}\quad\text{et}\quad B=\begin{pmatrix}
	B_{1,1}&\cdots&B_{1,q}\\
	\vdots&&\vdots\\
	B_{p,1}&\cdots&B_{p,q}
	\end{pmatrix}$$ avec des blocs \textit{compatibles pour le produit matriciel}, alors :
	$$AB=\begin{pmatrix}
		C_{1,1}&\cdots&C_{1,q}\\
		\vdots&&\vdots\\
		C_{n,1}&\cdots&C_{n,q}
	\end{pmatrix}\quad\text{avec}\quad\forall (i,j)\in\llbracket1,n\rrbracket\times\llbracket1,q\rrbracket,\ C_{i,j}=\sum_{k=1}^{p}A_{i,k}B_{k,q}.$$
\end{prop}
\begin{proof}
	Pour $(i,k,j)\in\llbracket1,n\rrbracket\times\llbracket1,p\rrbracket\times\llbracket1,q\rrbracket$, notons :
	\begin{itemize}
		\item $r_i$ le nombre de lignes des blocs de type $A_{i,*}$ ;
		\item $s_k$ le nombre de colonnes des blocs de type $A_{*,k}$ (qui est aussi le nombre de lignes des blocs de type $B_{k,*}$ car les blocs sont de tailles compatibles pour le produit) ;
		\item $t_j$ le nombre de colonnes des blocs de type $B_{*,j}$.
	\end{itemize}
	Écrivons la matrice $AB$ par blocs :
	$$AB=\begin{pmatrix}
		C_{1,1}&\cdots&C_{1,q}\\
		\vdots&&\vdots\\
		C_{n,1}&\cdots&C_{n,q}
	\end{pmatrix}\quad\text{avec}\quad\forall (i,j)\in\llbracket1,n\rrbracket\times\llbracket1,q\rrbracket,\ C_{i,j}\in\mathcal{M}_{r_i,t_j}(\K).$$
	Fixons alors $(i,j)\in\llbracket1,n\rrbracket\times\llbracket1,q$ et vérifions que $C_{i,j}=\displaystyle\sum_{k=1}^{p}A_{i,k}B_{k,q}$. \\
	Pour cela, fixons $(a,b)\in\llbracket1,r_i\rrbracket\times\llbracket1,t_j\rrbracket$ et montrons que 
	$$[C_{i,j}]_{a,b}=\left[\sum_{k=1}^{p}A_{i,k}B_{k,j}\right]_{a,b}.$$
	Notons $\alpha=\displaystyle\sum_{h=1}^{i-1}r_h$, $\gamma=\displaystyle\sum_{h=1}^{j-1}t_h$ et, pour $k\in\llbracket1,p+1\rrbracket$, $\beta_k=\displaystyle\sum_{h=1}^{k-1}s_h$. Remarquons que $\beta_{p+1}$ désigne le nombre de colonnes de $A$ et de lignes de $B$. On a d'une part :
	$$[C_{i,j}]_{a,b}=[AB]_{\alpha+a,\gamma+b}=\sum_{m=1}^{\beta_{p+1}}[A]_{\alpha+a,m}[B]_{m,\gamma+b}$$ et d'autre part :
	\begin{align*}
		\left[\sum_{k=1}^{p}A_{i,k}B_{k,j}\right]_{a,b}=\sum_{k=1}^{p}[A_{i,k}B_{k,j}]_{a,b}&=\sum_{k=1}^{p}\sum_{m=1}^{s_k}[A_{i,k}]_{a,m}[B_{k,j}]_{m,b}\\&=\sum_{k=1}^{p}\sum_{m=1}^{s_k}[A]_{\alpha+a,\beta_k+m}[B]_{\beta_k+m,\gamma+b}.
	\end{align*}
	Puisque les ensembles $(\llbracket\beta_k+1,\beta_k+s_k\rrbracket)_{k\in\llbracket1,p\rrbracket}$ forment un recouvrement disjoint de $\llbracket1,\beta_{p+1}\rrbracket$, on obtient finalement :
	$$\left[\sum_{k=1}^{p}A_{i,k}B_{k,j}\right]_{a,b}=\sum_{m=1}^{\beta_{p+1}}[A]_{\alpha+a,m}[B]_{m,\gamma+b}=[C_{i,j}]_{a,b},$$
	ce qui donne le résultat souhaité.
\end{proof}

\begin{exo}
	Soit $A\in\mathcal{M}_n(\K)$ et $M=\begin{pmatrix}
		0&A\\
		I_n&0
	\end{pmatrix}\in\mathcal{M}_{2n}(\K)$.
	\begin{enumerate}
		\item Expliciter les puissances de $M$.
		\item \begin{enumerate}[label=(\alph*)]
			\item Montrer que $\rg(M)=n+\rg(A)$.
			\item Montrer que $M$ est inversible si, et seulement si, $A$ l'est.
			\item Déterminer alors l'inverse de $A$.
		\end{enumerate}
	\end{enumerate}
\end{exo}

\begin{exo}
	\begin{enumerate}
		\item \begin{enumerate}[label=(\alph*)]
			\item Déterminer l'ensemble des matrices de $\mathcal{M}_n(\K)$ commutant avec la matrice $J_r\in\mathcal{M}_n(\K)$.
			\item Justifier qu'il s'agit d'un sous-espace vectoriel de $\mathcal{M}_n(\K)$ et donner sa dimension.
		\end{enumerate}
		\item Soit $p$ un projecteur de rang $r$ d'un $\K$-espace vectoriel $E$ de dimension $n$. Déterminer la dimension du sous-espace des endomorphismes de $E$ commutant avec $p$.
	\end{enumerate}
\end{exo}

\begin{attention}
	Le produit par blocs s'effectue comme un produit usuel mais, le produit matriciel n'étant pas commutatif, il est impératif de faire attention à l'ordre dans lequel on écrit chaque produit $A_{i,k}B_{k,j}$.
\end{attention}

\subsection{Matrices triangulaires/diagonales par blocs}

\begin{terminologie}
	On appelle matrice \textbf{triangulaire supérieure par blocs}, de blocs diagonaux $A_{1,1},\ldots,A_{n,n}$, une matrice carrée s'écrivant sous la forme :
	$$\begin{pmatrix}
		A_{1,1}&\cdots&\cdots&A_{1,n}\\
		0&\ddots&&\vdots\\
		\vdots&\ddots&\ddots&\vdots\\
		0&\cdots&0&A_{n,n}
	\end{pmatrix}$$ où les \textbf{blocs diagonaux} $A_{i,i}$ sont carrés.
	
	Sur le même principe, on parle de matrice \textbf{triangulaire inférieure par blocs} et de matrice \textbf{diagonale par blocs}.
\end{terminologie}

\begin{rmq}
	\begin{itemize}
		\item Couramment, dans l'écriture triangulaire par blocs, seuls les blocs diagonaux nous importent. On notera $\begin{pmatrix}
			A_1&&(\star)\\
			&\ddots&\\
			(0)&&A_n
		\end{pmatrix}$ une telle matrice triangulaire supérieure par blocs de blocs diagonaux $A_1,\ldots,A_n$.
		\item Un produit de deux matrices triangulaires supérieures par blocs, avec des blocs diagonaux de taille compatibles, l'est encore. Plus précisément :
		$$\begin{pmatrix}
			A_1&&(\star)\\
			&\ddots&\\
			(0)&&A_n
		\end{pmatrix}\begin{pmatrix}
		B_1&&(\star)\\
		&\ddots&\\
		(0)&&B_n
		\end{pmatrix}=\begin{pmatrix}
		A_1B_1&&(\star)\\
		&\ddots&\\
		(0)&&A_nB_n
		\end{pmatrix}.$$
		On a le même résultat pour des matrices triangulaires inférieures par blocs, et, évidemment, pour des matrices diagonales par blocs.
	\end{itemize}
\end{rmq}

\begin{notation}
	Une matrice diagonale par blocs de blocs diagonaux $A_1,\ldots,A_n$ est notée $\text{Diag}(A_1,\ldots,A_n)$. Ainsi, sous condition que les tailles des blocs soient compatibles, on a :
	$$\text{Diag}(A_1,\ldots,A_n)\text{Diag}(B_1,\ldots,B_n)=\text{Diag}(A_1B_1,\ldots,A_nB_n).$$
\end{notation}

\begin{met}
	Multiplier une matrice par une matrice diagonale par blocs $\text{Diag}(A_1,\ldots,A_n)$ :
	\begin{itemize}
		\item à gauche revient à multiplier les lignes par les blocs diagonaux $A_i$ ;
		\item à droite revient à multiplier les colonnes par les blocs diagonaux $A_j$.
	\end{itemize}
\end{met}

\begin{ex}
	\begin{enumerate}
		\item Pour tout $(A,B)\in\mathcal{M}_n(\K)\times\mathcal{M}_p(\K)$, on a $\rg\begin{pmatrix}
			A&0\\
			0&B
		\end{pmatrix}=\rg(A)+\rg(B)$.
		\item Plus généralement, le rang d'une matrice diagonale par blocs est égal à la somme des rangs des blocs diagonaux.\\
		Une matrice $A$ diagonale par blocs est donc inversible si, et seulement si, tous ses blocs diagonaux le sont. Dans ce cas, en notant $A_1,\ldots,A_r$ les blocs diagonaux de $A$, alors la matrice inverse $A\inv$ est diagonale par blocs de blocs diagonaux $A_1\inv,\ldots,A_r\inv$.
	\end{enumerate}
\end{ex}

\subsection{Transvections par blocs}

Dans ce qui suit, les tailles des matrices ne sont pas précisées (en particulier, on notera $I$ la matrice identité, sans préciser sa taille). Il est implicite qu'elles sont telles que les produits par blocs envisagés aient un sens.

\begin{itemize}
	\item Soit $M=\begin{pmatrix}
		A_1&B_1\\
		A_2&B_2
	\end{pmatrix}$ une matrice écrite par blocs. Lorsqu'il peut s'effectuer, on a le produit par blocs suivant :
	$$\begin{pmatrix}
		A_1&B_1\\
		A_2&B_2
	\end{pmatrix}\begin{pmatrix}
	I&0\\
	T&I
	\end{pmatrix}=\begin{pmatrix}
	A_1+B_1T&B_1\\
	A_2+B_2T&B_2
	\end{pmatrix}.$$
	On dit alors qu'on a effectué une \textbf{transvection par blocs}.
	
	\begin{notation}
		On notera $C\leftarrow C_1+C_2T$ la transvection par blocs ci-dessus.
	\end{notation}
	
	On peut, si l'on veut, "coder" la transvection par blocs avec des transvections "classiques" : si l'on a $T=(t_{i,j})_{(i,j)\in\llbracket1,n\rrbracket\times\llbracket1,p\rrbracket}$, alors la matrice $\begin{pmatrix}
		A_1+B_1T&B_1\\
		A_2+B_2T&B_2
	\end{pmatrix}$ est obtenue, à partir de $M$, appliquant les transvections suivantes (dont on voit aisément qu'elles commutent) :
	$$C_i\leftarrow C_i+t_{j,i}C_{j+n}\quad\text{avec}\quad(i,j)\in\llbracket1,n\rrbracket\times\llbracket1,p\rrbracket.$$
	La lourdeur de cette chaîne de transvections met en exergue l'efficacité et la concision apportée par l'utilisation des blocs.
	\item Avec les mêmes notations, on peut bien sûr modifier les blocs de droite à l'aide des blocs de gauche, en effectuant la transvection par blocs $C_2\leftarrow C_2+C_1T$ :
	$$\begin{pmatrix}
		A_1&B_1\\
		A_2&B_2
	\end{pmatrix}\begin{pmatrix}
		I&T\\
		0&I
	\end{pmatrix}=\begin{pmatrix}
		A_1&B_1+A_1T&\\
		A_2&B_2+A_2T
	\end{pmatrix}.$$
	\item On peut également réaliser des transvections par blocs en raisonnant sur les lignes (les deux transvections par blocs suivantes sont respectivement codées $L_1\leftarrow L_1+TL_2$ et $L_2\leftarrow L_2+TL_1$) :
	$$\begin{pmatrix}
		I&T\\
		0&I
	\end{pmatrix}\begin{pmatrix}
		A_1&A_2\\
		B_1&B_2
	\end{pmatrix}=\begin{pmatrix}
		A_1+TB_2&A_2+TB_2&\\
		B_1&B_2
	\end{pmatrix}$$
	et
	$$\begin{pmatrix}
		I&0\\
		T&I
	\end{pmatrix}\begin{pmatrix}
		A_1&A_2\\
		B_1&B_2
	\end{pmatrix}=\begin{pmatrix}
		A_1&A_2&\\
		B_1+TA_1&B_2+TA_2
	\end{pmatrix}.$$
\end{itemize}

\begin{rmq}
	Les matrices $\begin{pmatrix}
		I&T\\
		0&I
	\end{pmatrix}$ et $\begin{pmatrix}
	I&0\\
	T&I
	\end{pmatrix}$ sont triangulaires à diagonale composée uniquement de $1$, donc son inversibles et de déterminant $1$. Par conséquent, effectuer une transvection par blocs ne modifie ni le rang, ni le déterminant.
\end{rmq}

\begin{ex}
	Soit $(A,C)\in\mathcal{GL}_n(\K)\times\mathcal{M}_p(\K)$. Montrons que $\rg\begin{pmatrix}
		A&B\\
		0&C
	\end{pmatrix}=\rg(A)+\rg(C)$.
\end{ex}

Exercice Moulin.Mat.28

\begin{exo}
	Soit $A\in\mathcal{M}_{p,n}(\K)$ et $B\in\mathcal{M}(n,p)(\K)$. Montrer que :
	$$\rg\begin{pmatrix}
		I_n&B\\
		A&I_p
	\end{pmatrix}=\rg(I_p-AB)+n=\rg(I_n-BA)+p.$$
\end{exo}

\subsection{Stabilité et écriture par blocs}

Dans cette partie, $E$ est un espace vectoriel de dimension finie.

\begin{prop}[Traduction matricielle de la stabilité]%prop
	Soit $\mathcal{B}=(e_1,\ldots,e_n)$ une base de $E$ et $u\in\mathcal{L}(E)$. Si l'on écrit :
	$$M_\mathcal{B}(u)=\begin{pmatrix}
		A&C\\
		B&D
	\end{pmatrix}\quad\text{avec}\quad A\in\mathcal{M}_r(\K)$$
	alors :
	\begin{itemize}
		\item $\Vect(e_1,\ldots,e_r)$ est stable par $u$ si, et seulement si, $B=0$ ;
		\item $\Vect(e_{r+1},\ldots,e_n)$ est stable par $u$ si, et seulement si, $C=0$.
	\end{itemize}
\end{prop}
\begin{proof}
	Notons $F=\Vect(e_1,\ldots,e_r)$.\\
	La stabilité de $F$ se traduit par le fait que les $r$ première colonnes de sa matrice dans $\mathcal{B}$ représentent, dans $\mathcal{B}$, des vecteurs de $F$, c'est-à-dire  que ces colonnes aient leurs $n-r$ derniers coefficients non nuls. Cela signifie exactement que $B=0$. \\
	Démonstration analogue pour la stabilité de $\Vect(e_{r+1},\ldots,e_n)$.
\end{proof}

\begin{rmq}
	Reprenons les notations précédentes et notons :
	$$F=\Vect(e_1,\ldots,e_r)\quad\text{et}\quad G=\Vect(e_{r+1},\dots,e_n).$$
	\begin{itemize}
		\item Dans le cas où $B=0$ :
		\begin{enumerate}[label=$\star$]
			\item $A$ est la matrice dans la base $(e_1,\ldots,e_r)$ de $F$ de l'endomorphisme induit $u_F$ ;
			\item l'interprétation de $D$ est plus délicate : en notant $p\in\mathcal{L}(E)$ la projection sur $G$ parallèlement à $F$, alors $D$ est la matrice dans la base $(e_{r+1},\ldots,e_n)$ de $G$ de l'endomorphisme induit par $p\circ u$ sur $G$.
		\end{enumerate}
		\item De même, dans le cas où $C=0$, $D$ est la matrice dans la base $(e_{r+1},\ldots,e_n)$ de l'endomorphisme induit par $u$ sur $G$.
	\end{itemize}
\end{rmq}

\begin{ex}
	Soit $\mathcal{B}=(e_1,\ldots,e_n)$ une base de $E$ et $u\in\mathcal{L}(E)$.
	\begin{itemize}
		\item La matrice de $u$ dans la base $\mathcal{B}$ est triangulaire supérieure si, et seulement si :
		$$\forall i\in\llbracket1,n\rrbracket,\ u(e_i)\in\Vect(e_1,\ldots,e_i),$$ c'est-à-dire si, et seulement si, $u$ laisse stable chacun des $F_i=\Vect(e_1,\ldots,e_i)$.
		\item On retrouve le fait que si $M_\mathcal{B}(u)$ est triangulaire supérieure et si $u$ est bijectif, alors $M_\mathcal{B}(u\inv)$ est également triangulaire supérieure. En effet, si c'est le cas, alors, pour tout $i\in\llbracket1,n\rrbracket$, on a $u(F_i)\subset F_i$ puis, comme $u$ est bijectif, $\dim(u(F_i))=\dim(F_i)$ donc $u(F_i)=F_i$ et enfin $F_i=u\inv(u(F_i))=u\inv(F_i)$.
	\end{itemize}
\end{ex}

\begin{prop}%prop
	Soit $E$ de dimension finie tel que $E=\displaystyle\bigoplus_{i=1}^pE_i$. Si $\mathcal{B}$ est une base adaptée à cette décomposition, alors l'endomorphisme $u$ stabilise chaque $E_i$ si, et seulement si, sa matrice dans la base $\mathcal{B}$ est diagonale par blocs :
	$$\begin{pmatrix}
		A_1&&(0)\\
		&\cdots&\\
		(0)&&A_p
	\end{pmatrix},$$ où, pour tout $i\in\llbracket1,p\rrbracket$, la matrice $A_i$ est carrée de taille $\dim(E_i)$.\\
	Dans ce cas, chacun des blocs diagonaux $A_i$ est la matrice de l'endomorphisme induit par $u$ sur $E_i$ dans la base de $E_i$ extraite de $\mathcal{B}$.
\end{prop}
\begin{proof}
	Pour $i\in\llbracket1,p\rrbracket$, notons $\mathcal{B}_i$ la base de $E_i$ extraite de $\mathcal{B}$ ; le sous-espace $E_i$ est alors stable par $u$ si, et seulement si, les vecteurs de $\mathcal{B}_i$ ont pour image par $u$ un vecteur de $E_i$, c'est-à-dire une combinaison linéaire de vecteurs de $\mathcal{B}_i$. D'où l'équivalence annoncée et l'interprétation des blocs diagonaux $A_i$.
\end{proof}

\begin{ex}
	Soit $\mathcal{B}=(e_1,\ldots,e_n)$ une base de $E$, et $u\in\mathcal{L}(E)$. La matrice de $u$ dans la base $\mathcal{B}$ est diagonale si, et seulement si, $u$ stabilise chaque droite $\Vect(e_i)$ pour $i\in\llbracket1,n\rrbracket$.
\end{ex}

\section{Déterminants}

\subsection{Formes $n$-linéaires alternées}

\begin{dfn}
	Une forme $n$-linéaire $f$ sur $E$ est dite \textbf{alternée} si, pour tous vecteurs $u_1,\ldots,u_n$ de $E$ et pour tout $(i,j)\in\llbracket1,n\rrbracket^2$ avec $i\ne j$, on a :
	$$u_i=u_j\implies f(u_1,u_2,\ldots,u_n)=0.$$
\end{dfn}

\begin{prop}%prop
	Si $f$ est une forme $n$-linéaire alternée, alors elle est \textbf{antisymétrique}, c'est-à-dire qu'elle vérifier, pour tout $(u_1,\ldots,u_n)\in E^n$ et pour tout $(i,j)\in\llbracket1,n\rrbracket^2$ tel que $i<j$ :
	$$f(u_1,\ldots,\underset{\substack{\uparrow\\i\text{-ème}}}{u_i},\ldots,\underset{\substack{\uparrow\\j\text{-ème}}}{u_j},\ldots,u_n)=-f(u_1,\ldots,\underset{\substack{\uparrow\\i\text{-ème}}}{u_j},\ldots,\underset{\substack{\uparrow\\j\text{-ème}}}{u_i},\ldots,u_n).$$
\end{prop}

\begin{exo}
	Montrer que la réciproque de la proposition précédente est vrai si $\K$ est un corps de caractéristique différente de $2$. Comme nous travaillons avec des sur-corps de $\Q$, cela sera donc toujours vrai.
\end{exo}

\begin{thm}%thm
	Soit $(e_1,\ldots,e_n)$ une base de $E$.
	\begin{itemize}
		\item Il existe une unique forme $n$-linéaire alternée $\varphi_0$ sur $E$ telle que :
		$$\varphi_0(e_1,\ldots,e_n)=1.$$
		\item Toute forme $n$-linéaire alternée sur $E$ est proportionnelle à $\varphi_0$.
	\end{itemize}
\end{thm}

Exercice Moulin.Mat.53

\subsection{Déterminant d'une famille de vecteurs}

\begin{dfn}
	Soit $\mathcal{B}=(e_1,\ldots,e_n)$ une base de $E$. Le \textbf{déterminant dans la base $\mathcal{B}$}, noté $\det_\mathcal{B}$, est l'unique forme linéaire alternée sur $E$ telle que $\det_\mathcal{B}(\mathcal{B})=1$.\\
	Si $u_1,\ldots,u_n$ sont des vecteurs de $E$ et $(a_{i,j})_{1\leqslant i,j\leqslant n}$ des scalaires tels que :
	$$\forall j\in\llbracket1,n\rrbracket,\ u_j=\sum_{i=1}^{n}a_{i,j}e_i,$$ alors on a :
	$$\det_\mathcal{B}(u_1,\ldots,u_n)=\sum_{\sigma\in\mathcal{S}_n}\varepsilon(\sigma)a_{\sigma(1),1}\cdots a_{\sigma(n),n}.$$
\end{dfn}

\begin{exo}
	Rappeler les propriétés de ce déterminant.
\end{exo}

\begin{exo}
	Quel est le coefficient de proportionnalité entre $\det_\mathcal{B}$ et $\det_{\mathcal{B}'}$, si $\mathcal{B}$ et $\mathcal{B}'$ sont deux bases de $E$ ?
\end{exo}

\subsection{Déterminant d'un endomorphisme}

Soit $E$ un espace vectoriel de dimension finie $n\geqslant 1$.

\begin{dfn}
	Soit $f$ un endomorphisme de $E$. Il existe alors un unique scalaire $\lambda$, appelé \textbf{déterminant} de $f$, tel que pour toute base $\mathcal{B}$ de $E$ et, pour tout $(u_1,\ldots,u_n)\in E^n$, on ait :
	$$\det_\mathcal{B}(f(u_1),\ldots,f(u_n))=\lambda\det_\mathcal{B}(u_1,\ldots,u_n).$$
	Le déterminant de $f$ se note $\det(f)$ ou $\det f$ et, pour toute base $\mathcal{B}$ de $E$, on a :
	$$\det(f)=\det_\mathcal{B}(f(\mathcal{B})).$$
\end{dfn}

\begin{exo}
	Rappeler les propriétés de ce déterminant.
\end{exo}
% det(u\circ v)=det(u)det(v)

\begin{exo}
	Soit $E$ et $F$ deux espaces vectoriels normés de dimension finie. \\
	A-t-on $\det(v\circ u)=\det(u\circ v)$ pour tout $u\in\mathcal{L}(E,F)$ et tout $v\in\mathcal{L}(F,E)$ ?
\end{exo}

\subsection{Déterminant d'une matrice carrée}

\begin{dfn}
	On appelle \textbf{déterminant d'une matrice carrée} d'ordre $n$ le déterminant de ses vecteurs colonnes dans la base canonique de $\K^n$. On le note $\det(A)$ ou $\det A$. \\
	Si $A=(a_{i,j})_{1\leqslant i,j\leqslant n}$ est une matrice carrée d'ordre $n$, alors on a :
	$$\det(A)=\sum_{\sigma\in\mathcal{S}_n}\varepsilon(\sigma)\prod_{i=1}^{n}a_{\sigma(i),i}.$$
\end{dfn}

\begin{exo}
	Rappeler les propriétés de ce déterminant.
\end{exo}
%det(u_A)=det(A), det(A^T)=det(A)

Exercice Moulin.Mat.56

\subsection{Application : orientation d'un $\R$-espace vectoriel}

Soit $E$ un $\R$-espace vectoriel de dimension finie $n\geqslant 1$.\\
On définit sur l'ensemble des bases de $E$ la relation suivante :
$$\mathcal{B}\mathcal{R}\mathcal{B}'\iff\det_\mathcal{B}(\mathcal{B}')>0.$$
Lorsque $\mathcal{B}\mathcal{R}\mathcal{B}'$, on dit que les bases $\mathcal{B}$ et $\mathcal{B}'$ ont \textbf{même orientation}.

\begin{prop}%prop
	La relation $\mathcal{R}$ est une relation d'équivalence sur l'ensemble des bases de $E$. Elle possède deux classes d'équivalence.
\end{prop}
\begin{proof}
	Les propriétés usuelles du déterminant montrent aisément qu'il s'agit bien d'une relation d'équivalence. Pour les classes d'équivalence, en fixant $\mathcal{B}_0=(e_1,\ldots,e_n)$ une base de $E$ et en posant $\mathcal{B}_0'=(-e_1,e_2,\ldots,e_n)$, les deux classes d'équivalence sont les classes de $\mathcal{B}_0$ et de $\mathcal{B}_0'$.
\end{proof}

Choisir une orientation de $E$, c'est choisir une classe d'équivalence et donc un représentant $\mathcal{B}_0$ de cette classe. Fixons une telle base.

\begin{dfn}
	On dit qu'une base $\mathcal{B}$ de $E$ est :
	\begin{itemize}
		\item \textbf{directe} si $\mathcal{B}\mathcal{R}\mathcal{B}_0$, c'est-à-dire si $\det_{\mathcal{B}_0}(\mathcal{B})>0$ ;
		\item \textbf{indirecte} sinon.
	\end{itemize}
\end{dfn}

\begin{rmq}
	En pratique, dans $\R^n$, on choisit toujours la base canonique comme base directe de référence.
\end{rmq}

\begin{attention}
	Orienter un espace vectoriel ne permet pas d'orienter naturellement ses sous-espaces vectoriels.
	%Exemple de l'espace et du plan du tableau
	%On peut orienter un hyperplan en complétant la base en une base directe de l'espace.
\end{attention}

\subsection{Calcul des déterminants}

Revoir dans le cours de première année les méthodes usuelles de calcul de déterminant :
\begin{itemize}
	\item opérations élémentaires ;
	\item développement par rapport à une rangée. %(-1)^{i+J}$, préférer la distance à la diagonale
\end{itemize}

\subsection{Déterminant et écriture par blocs}

\paragraph{Convention} Si une matrice est de taille nulle, son déterminant vaut $1$.

\begin{prop}%prop
	Le déterminant d'une matrice triangulaire par blocs est le produit des déterminants de ses blocs diagonaux, c'est-à-dire :
	$$\det\begin{pmatrix}
		A_1&&(\star)\\
		&\ddots&\\
		(0)&&A_n
	\end{pmatrix}=\det\begin{pmatrix}
	A_1&&(0)\\
	&\ddots&\\
	(\star)&&A_n
	\end{pmatrix}=\prod_{k=1}^{n}\det(A_k).$$
\end{prop}

\begin{proof}
	Quitte à faire une récurrence sur le nombre de blocs, il suffit de prouver le résultat pour une matrice $L\in\mathcal{M}_n(\K)$ de la forme $M=\begin{pmatrix}
		A&C\\
		0&B
	\end{pmatrix}$.\\
	Pour cela, on remarque que $M=\begin{pmatrix}
		I_r&C\\
		0&B
	\end{pmatrix}\begin{pmatrix}
	A&0\\
	0&I_{n-r}
	\end{pmatrix}$.
\end{proof}

Exercice Moulin.Mat.71 en utilisant l'exercice Moulin.Mat.56.

Exercice Moulin.Mat.70.

\begin{rmq}
	Une matrice triangulaire par blocs est donc inversible si, et seulement si, tous ses blocs diagonaux sont inversibles.
\end{rmq}

\begin{rmq}
	\begin{enumerate}
		\item On a, en reconnaissant un déterminant triangulaire par blocs :
		$$\begin{vmatrix}
			1&2&31&45\\
			3&4&10&23\\
			0&0&1&3\\
			0&0&1&1
		\end{vmatrix}=\begin{vmatrix}
		1&2\\
		3&4
		\end{vmatrix}\times\begin{vmatrix}
		1&3\\
		1&1
		\end{vmatrix}=(-2)\times(-2)=4.$$
		\item Soit $(A,B)\in\mathcal{M}_n(\K)^2$. Montrons que $\det\begin{pmatrix}
			A&B\\
			B&A
		\end{pmatrix}=\det(A+B)\det(A-B)$.
	\end{enumerate}
\end{rmq}

\subsection{Déterminant de Vandermonde}

\begin{prop}%prop
	Étant donnés des scalaires $x_0,x_1,\ldots,x_n$, on note $V(x_0,x_1,\ldots,x_n)$ le déterminant d'ordre $n+1$ défini par :
	$$V(x_0,x_1,\ldots,x_n)=\begin{vmatrix}
		1&1&\cdots&1\\
		x_0&x_1&\cdots&x_n\\
		\vdots&\vdots&&\vdots\\
		x_0^n&x_1^n&\cdots&x_n^n
	\end{vmatrix}.$$
	\begin{enumerate}
		\item $V(x_0,\ldots,x_n)$ est non nul si, et seulement si, les $x_i$ sont distincts deux à deux.
		\item Dans tous les cas, on a :
		$$V(x_0,\ldots,x_n)=\prod_{0\leqslant j<i\leqslant n}(x_i-x_j).$$
	\end{enumerate}
\end{prop}

\begin{exo}[Vandermonde et polynômes]
	Soit $A=\begin{pmatrix}
		1&x_0&\cdots&x_0^n\\
		1&x_1&\cdots&x_1^n\\
		\vdots&\vdots&&\vdots\\
		1&x_n&\cdots&x_n^n
	\end{pmatrix}$.
	\begin{enumerate}
		\item Redémontrer directement le premier résultat de la proposition précédente en considérant le système homogène associé à la matrice $A$.
		\item Lorsque $A$ est inversible, expliciter son inverse.
	\end{enumerate}
\end{exo}

\subsection{Comatrice}

\begin{dfn}
	Étant donnée une matrice $A=(a_{i,j})_{\substack{1\leqslant i\leqslant n\\
	1\leqslant j\leqslant n}}\in\mathcal{M}_n(\K)$, on appelle \textbf{comatrice} de $A$, notée $\text{Com}(A)$ ou $\text{Com}A$, la matrice des cofacteurs de $A$, c'est-à-dire la matrice $B=(b_{i,j})_{\substack{1\leqslant i\leqslant n\\
	1\leqslant j\leqslant n}}\in\mathcal{M}_n(\K)$, où $b_{i,j}$ est le cofacteur de $a_{i,j}$ dans $A$.
\end{dfn}

\begin{prop}%prop
	Si $B$ est la comatrice de $A\in\mathcal{M}_n(\K)$, on a :
	$$AB^T=B^TA=(\det(A))I_n.$$
\end{prop}

Exercice Moulin.Mat.78

\begin{cor}%cor
	Si $B$ est la comatrice d'une matrice $A\in\mathcal{GL}_n(\K)$, alors on a :
	$$A\inv=\frac{1}{\det(A)}B^T.$$
\end{cor}

\begin{met}
	On utilise cette formule pour inverser pratiquement une matrice $n\times n$ si, et seulement si, $n=2$.
\end{met}

\begin{exo}
	Déterminer les éléments inversibles de l'anneau $\mathcal{M}_n(\Z)$.
\end{exo}

\begin{exo}
	Soit $A:t\mapsto A(t)$ une application continue à valeurs dans $\mathcal{GL}_n(\K)$. Montrer que $t\mapsto A(t)\inv$ est continue.
\end{exo}

\subsection{Formules de Cramer (HP)}

Considérons un système de Cramer $AX=B$, où $A$ est une matrice inversible de $\mathcal{M}_n(\K)$ et $B\in\K^n$. Pour tout $k\in\llbracket1,n\rrbracket$, on note $A_k$ la matrice obtenue en remplaçant la $k$-ième colonne de $A$ par $B$, le seconde membre du système.

\begin{thm}[Formules de Cramer]%thm
	L'unique solution du système est donnée par :
	$$\forall k\in\llbracket1,n\rrbracket,\ x_k=\frac{\det(A_k)}{\det(A)}.$$
\end{thm}
\begin{proof}
	Décomposer $B$ dans la base de $\K^n$ constituée des colonnes de $A$, et utiliser la multilinéarité du déterminant.
\end{proof}

\begin{met}
	On utilise cette formule pour résoudre pratiquement un système $n\times n$ si, et seulement si, $n=2$.
\end{met}

\begin{exo}
	Soit $(O(t),\overline{u}(t),\overline{v}(t),\overline{w}(t))$ un repère mobile de l'espace, et $M(t)$ un point mobile du plan $(O(t),\overline{u}(t),\overline{v}(t))$. On suppose que les fonctions $O$, $M$, $\overline{u}$, $\overline{v}$ et $\overline{w}$ sont continues.\\
	Montrer que les coordonnées de $M(t)$ dans $(O(t),\overline{u}(t),\overline{v}(t))$ sont continues.
\end{exo}

\end{document}